% Mathématiques 3e — V3 LaTeX
% Proportionnalité, taux d’évolution et coefficient multiplicateur

\section{Objectifs}
\begin{itemize}
\item Relier une situation proportionnelle à une fonction linéaire.
\item Calculer et interpréter un taux d’évolution.
\item Passer d’un taux d’évolution à un coefficient multiplicateur.
\item Résoudre des problèmes d’échelles et d’évolutions successives.
\end{itemize}

\section{Repères et vocabulaire}
En troisième, la proportionnalité devient un pont entre calcul numérique, pourcentages, fonctions et géométrie. Le taux d’évolution et le coefficient multiplicateur permettent de décrire précisément une hausse ou une baisse et de comparer des évolutions successives.

\section{1. Reconnaître la proportionnalité}
Une relation est proportionnelle lorsque le quotient entre deux grandeurs correspondantes reste constant. Elle se modélise par $y=kx$.

\[
y=3{,}5x
\]

\begin{exemple}
Graphiquement, une situation proportionnelle est représentée par une droite passant par l’origine.
\end{exemple}

\section{2. Taux d’évolution}
Le taux d’évolution mesure la variation relative entre une valeur initiale et une valeur finale.

\[
t=\dfrac{V_f-V_i}{V_i}
\]

\begin{exemple}
Un taux positif décrit une hausse ; un taux négatif décrit une baisse.
\end{exemple}

\coursefigure{/images/mathematiques/3eme/proportionnalite-taux-evolution-1.svg}{Représentation structurante pour le chapitre Proportionnalité, taux d’évolution et coefficient multiplicateur.}{Cette figure organise visuellement les relations utilisées dans les premières notions du chapitre.}

\section{3. Coefficient multiplicateur}
Une évolution de taux $t$ transforme une valeur $V_i$ en $V_f=V_i(1+t)$. Le nombre $1+t$ est le coefficient multiplicateur.

\[
+12\,\%\Rightarrow k=1{,}12
\]

\[
-20\,\%\Rightarrow k=0{,}80
\]

\begin{exemple}
Le coefficient multiplicateur évite de confondre la variation avec la valeur finale.
\end{exemple}

\section{4. Retrouver un taux}
Lorsque le coefficient multiplicateur est connu, le taux vaut $k-1$.

\[
k=1{,}075\Rightarrow t=0{,}075=7{,}5\,\%
\]

\begin{exemple}
Une valeur finale plus petite que la valeur initiale conduit à un coefficient inférieur à $1$.
\end{exemple}

\section{5. Évolutions successives}
Deux évolutions successives se composent en multipliant leurs coefficients multiplicateurs, pas en additionnant directement les pourcentages.

\[
1{,}10\times0{,}90=0{,}99
\]

\begin{exemple}
Une hausse de $10\,\%$ suivie d’une baisse de $10\,\%$ produit une baisse globale de $1\,\%$.
\end{exemple}

\coursefigure{/images/mathematiques/3eme/proportionnalite-taux-evolution-2.svg}{Deuxième représentation pédagogique pour le chapitre Proportionnalité, taux d’évolution et coefficient multiplicateur.}{Cette représentation sert de support de lecture, de comparaison ou de contrôle dans les problèmes du chapitre.}

\section{6. Échelle}
Une échelle est une relation proportionnelle entre une représentation et la réalité. Il faut rendre les unités compatibles avant d’appliquer le coefficient.

\[
1:500\Rightarrow k=\dfrac1{500}
\]

\begin{exemple}
Une longueur de $8{,}4\,\text{cm}$ sur le plan représente $42\,\text{m}$ à l’échelle $1:500$.
\end{exemple}

\section{7. Proportionnalité et géométrie}
Thalès, triangles semblables et homothéties utilisent des rapports de longueurs constants. La proportionnalité organise les correspondances géométriques.

\[
\dfrac{AM}{AB}=\dfrac{AN}{AC}
\]

\begin{exemple}
On ne peut utiliser un rapport que si les longueurs comparées sont bien correspondantes.
\end{exemple}

\section{8. Comparer deux modèles}
Un tarif peut être proportionnel tandis qu’un autre comporte un coût fixe. Les fonctions et inéquations permettent alors de déterminer un seuil.

\[
A(x)=12+4x
\]

\[
B(x)=5{,}5x
\]

\begin{exemple}
Le calcul du point d’égalité ne suffit pas : il faut aussi comparer avant et après ce seuil.
\end{exemple}

\section{Méthode générale de résolution}
\begin{methode}
\begin{enumerate}
\item Identifier les valeurs initiale et finale ou les grandeurs correspondantes.
\item Mettre les unités sous une forme compatible.
\item Choisir entre coefficient de proportionnalité, taux d’évolution et coefficient multiplicateur.
\item Écrire la relation algébrique avant d’effectuer le calcul.
\item Contrôler si le résultat décrit une hausse, une baisse ou une échelle plausible.
\item Interpréter le résultat dans la situation réelle.
\end{enumerate}
\end{methode}

\section{Problème guidé et stratégie}
Dans un problème de troisième, la difficulté ne vient pas seulement du calcul. Il faut reconnaître la notion utile, choisir une représentation, enchaîner plusieurs étapes et expliquer pourquoi la méthode est valide. On commence donc par isoler les données, puis on écrit la relation mathématique avant de remplacer par des nombres. Une réponse est considérée comme complète lorsqu’elle comporte une justification, un calcul lisible, un contrôle et une interprétation.

\begin{attention}
Une figure, un graphique, un tableau ou une calculatrice fournit des informations, mais ne remplace pas une justification lorsque le problème demande une preuve. Inversement, un calcul exact sans interprétation peut rester incomplet dans un contexte concret.
\end{attention}

\section{Contrôler un résultat}
Avant de valider une réponse, on vérifie le signe, l’ordre de grandeur, les unités et la compatibilité avec les hypothèses. On peut aussi refaire le calcul par une autre représentation : formule et graphique, valeur exacte et approximation, calcul direct et vérification dans l’énoncé. Cette habitude permet de repérer une erreur de saisie ou une mauvaise modélisation avant la conclusion.

\section{Erreurs fréquentes}
\begin{itemize}
\item Additionner des pourcentages successifs sans justification.
\item Confondre taux et coefficient multiplicateur.
\item Appliquer une échelle avec des unités différentes.
\item Reconnaître une proportionnalité sans vérifier le passage par l’origine.
\item Utiliser une quatrième proportionnelle dans une situation non proportionnelle.
\item Oublier d’interpréter le signe du taux.
\end{itemize}

\section{À retenir}
\begin{aretenir}
\begin{itemize}
\item Une proportionnalité se modélise par $y=kx$.
\item Le taux d’évolution vaut $(V_f-V_i)/V_i$.
\item Le coefficient multiplicateur vaut $1+t$.
\item Les évolutions successives se composent en multipliant les coefficients.
\end{itemize}
\end{aretenir}

\section{Réinvestissement : Proportionnalité et géométrie}
Pour réinvestir cette notion, on part d’une situation où plusieurs représentations sont possibles. Thalès, triangles semblables et homothéties utilisent des rapports de longueurs constants. La proportionnalité organise les correspondances géométriques. L’élève doit expliciter son choix, effectuer le calcul et revenir au contexte. On ne peut utiliser un rapport que si les longueurs comparées sont bien correspondantes. Le contrôle final consiste à vérifier que la réponse respecte les hypothèses et que son ordre de grandeur est compatible avec les données.

\section{Réinvestissement : Coefficient multiplicateur}
Pour réinvestir cette notion, on part d’une situation où plusieurs représentations sont possibles. Une évolution de taux $t$ transforme une valeur $V_i$ en $V_f=V_i(1+t)$. Le nombre $1+t$ est le coefficient multiplicateur. L’élève doit expliciter son choix, effectuer le calcul et revenir au contexte. Le coefficient multiplicateur évite de confondre la variation avec la valeur finale. Le contrôle final consiste à vérifier que la réponse respecte les hypothèses et que son ordre de grandeur est compatible avec les données.

\section{Réinvestissement : Taux d’évolution}
Pour réinvestir cette notion, on part d’une situation où plusieurs représentations sont possibles. Le taux d’évolution mesure la variation relative entre une valeur initiale et une valeur finale. L’élève doit expliciter son choix, effectuer le calcul et revenir au contexte. Un taux positif décrit une hausse ; un taux négatif décrit une baisse. Le contrôle final consiste à vérifier que la réponse respecte les hypothèses et que son ordre de grandeur est compatible avec les données.
